\section*{ВВЕДЕНИЕ}
\addcontentsline{toc}{section}{ВВЕДЕНИЕ}

Игровые движки играют ключевую роль в обеспечении разработки современных компьютерных игр и интерактивных трёхмерных приложений. В условиях стремительного роста требований к графике, производительности и кросс-платформенности разработка качественного игрового движка становится особенно важной. Однако традиционный подход к созданию игр на основе готовых коммерческих решений, таких как Unity или Unreal Engine, не всегда даёт полное понимание низкоуровневых механизмов рендеринга, управления ресурсами и архитектуры игровых систем. Отсутствие глубоких знаний в области работы с графическим API, шейдерами и управлением памятью существенно ограничивает возможности разработчика при создании специализированных игровых проектов.

Разработка собственного игрового движка позволяет с нуля реализовать все ключевые компоненты: графическую подсистему, систему управления игровыми объектами, менеджер сцен, систему камер и освещения, а также визуальный редактор сцен. Проведённый анализ существующих игровых движков (Unity, Unreal Engine, Godot) показал, что они либо обладают избыточным функционалом, не требуемым в учебных и небольших игровых проектах, либо предоставляют ограниченные возможности для изучения внутреннего устройства. В связи с этим целесообразной является разработка собственного программного комплекса с визуальной средой проектирования для создания компьютерных игр.

\emph{Цель настоящей работы} – разработка и внедрение игрового движка для создания 3D интерактивных приложений в реальном времени. Для достижения поставленной цели необходимо решить \emph{следующие задачи:}
\begin{itemize}
\item провести анализ предметной области;
\item разработать концептуальную модель движка;
\item спроектировать движок;
\item реализовать сайт средствами выбранных технологий (C\#, OpenTK, Windows Forms, SixLabors.ImageSharp).
\end{itemize}

\emph{Структура и объем работы.} Отчет состоит из введения, 4 разделов основной части, заключения, списка использованных источников, 2 приложений. Текст выпускной квалификационной работы равен \formbytotal{lastpage}{страниц}{е}{ам}{ам}.

\emph{Во введении} сформулирована цель работы, поставлены задачи разработки, описана структура работы, приведено краткое содержание каждого из разделов.

\emph{В первом разделе} на стадии описания технической характеристики предметной области приводится сравнение существующих решений. Сравниваются современные архитектурные решения. 

\emph{Во втором разделе} на стадии технического задания формируются требования к функционалу движка, программной и аппаратной совместимости, надежности, графическому API, производительности, безопасности, размещению и сопровождению и оформлению документации.

\emph{В третьем разделе} на стадии технического проектирования представлена архитектура разрабатываемого решения и обоснование выбора технологий.

\emph{В четвертом разделе} приводится список классов и их не публичных методов, использованных при разработке движка, производится его тестирование.

В заключении излагаются основные результаты работы, полученные в ходе разработки.

В приложении А представлен графический материал.
В приложении Б представлены фрагменты исходного кода. 
